\section{绪论}
\subsection{问题来源与研究目标}
教育研究应把一般教学关切转化为具体、可回答的问题。例如“如何识别课堂提问的层次与反馈方式”比“如何提高教育质量”具有更明确的观察边界。选题时还需确认研究进入现场的条件、时间安排与资料使用权限。
\subsection{概念界定与文献综述}
将课堂提问、候答时间、追问及反馈等概念作操作性定义。文献综述应比较已有研究的方法与证据，不能将理论分类直接当成效果结论。本例引用真实的教育分类和质性分析著作，演示来源标注\cite{anderson2001,braun2006}；它们不构成本选题的完整综述。
\subsubsection{研究边界}
明确学段、学科、课型与分析单位，说明哪些问题不在研究范围内。课堂记录包含未成年人信息时，应按机构要求取得必要审批、监护人同意和参与者适当同意，不应凭模板文字声称已经获批\footnote{本脚注演示补充说明，模板没有收集任何学生或教师资料。}。
\section{研究设计}
\subsection{抽样与资料来源}
说明学校、教师和课堂的选择方式以及纳入、排除标准。便利抽样不能自动代表总体，不同课堂的观察次数也不等于独立教师数量。正文只记录实际完成的工作，不预填访谈人数、课时数量或场所名称。
\subsection{研究流程}
图~\ref{fig:education} 给出一条可供修改的工作路线，不是学校要求的固定方法。
\begin{figure}[htbp]\centering
\begin{tikzpicture}[node distance=6mm,b/.style={draw,minimum width=65mm,minimum height=9mm,font=\small},>=Stealth]
\node[b](a){界定问题与取得必要许可};
\node[b,below=of a](b){试用观察工具并修订编码规则};
\node[b,below=of b](c){收集资料与匿名化整理};
\node[b,below=of c](d){分析、复核与报告局限};
\draw[->](a)--(b);\draw[->](b)--(c);\draw[->](c)--(d);
\end{tikzpicture}
\caption{研究组织流程示例}\label{fig:education}
\end{figure}
\subsection{观察维度与编码规则}
表~\ref{tab:coding} 只展示编码表结构。正式研究需通过试编码明确类别边界，处理难以判断的情况，并记录规则修改。
\begin{table}[htbp]\centering\caption{观察维度示例（非调查结果）}\label{tab:coding}
\begin{tabularx}{\linewidth}{lXX}\toprule 维度 & 记录内容 & 注意事项\\\midrule
问题类型 & 问题原文与分类依据 & 不以关键词机械判断认知水平\\
候答时间 & 提问结束至首次回应的间隔 & 明确起止点与计时规则\\
教师反馈 & 追问、转问、解释等行为 & 保留上下文，允许多重编码\\
参与方式 & 自愿、点名或集体回答 & 不公开学生身份\\\bottomrule\end{tabularx}
\end{table}
\subsection{分析方法与一致性}
若使用分类统计，应说明分析单位、分母和缺失处理。例如某类问题占比可写为：
\begin{equation}\label{eq:proportion}p_k=\frac{n_k}{N}\times100\%.\end{equation}
式中，$n_k$ 是第 $k$ 类事件数，$N$ 是满足统计条件的总事件数。事件数不等于独立样本量，重复课堂不能任意当作独立观测。质性分析应呈现编码形成过程和反例，不能只用少量摘录印证既定观点。
\section{结果报告的组织方式}
\subsection{描述性结果}
实际研究应先报告资料规模、缺失和排除情况，再呈现与问题对应的结果。本模板不提供真实数据，因此没有柱状图、比例或访谈结论可报告。图表应有正文引导，资料摘录应匿名并注明内部标识。
\subsection{证据与解释的区分}
分开呈现观察事实、受访者的观点和研究者的解释。单次课堂中的沉默可能有多种原因，不能直接等同于学生不会思考；受访者的自述也不等于课堂行为的客观记录。
\section{讨论与建议}
\subsection{与已有研究对话}
比较研究发现与文献中的一致和不一致之处，解释方法差异是否可能影响结论。建议应由证据支持，并说明资源、教师经验及课型等实施条件，避免把一般常识包装成独立研究发现。
\subsection{局限与研究伦理}
说明抽样范围、观察者效应、编码主观性及资料不完整等限制。保存录音、视频和原始记录时应控制访问权限，不能把个人信息随论文公开。最终提交前还应核对学院规定的论文格式与归档材料。
